\section{STARKs and the Algebraic Intermediate Representation (AIR)}

\subsection{Motivation}

Traditional zk-SNARK systems represent computation as static constraint systems such as R1CS or PLONK circuits.  
In contrast, STARKs describe computation as a dynamic process evolving over time.

\medskip

\noindent
\textbf{Key idea.}  
A computation is represented as a sequence of states, and correctness is enforced by algebraic constraints on state transitions.

\subsection{Execution Trace}

An \emph{execution trace} is a table:
\[
\begin{array}{c|ccc}
t & s_t^{(1)} & s_t^{(2)} & \cdots \\
\hline
0 & s_0^{(1)} & s_0^{(2)} & \cdots \\
1 & s_1^{(1)} & s_1^{(2)} & \cdots \\
2 & s_2^{(1)} & s_2^{(2)} & \cdots \\
\vdots & \vdots & \vdots & \ddots
\end{array}
\]

Each row represents the system state at time step $t$.

\subsection{AIR Constraints}

The Algebraic Intermediate Representation consists of:

\paragraph{Transition constraints.}
\[
P(s_t, s_{t+1}) = 0.
\]

Example:
\[
x_{t+1} - x_t^2 - 1 = 0.
\]

\paragraph{Boundary constraints.}
\[
s_0 = \text{input}, \qquad s_T = \text{output}.
\]

\paragraph{Algebraic form.}
All constraints must be polynomial equations over a finite field.

\subsection{Example}

Consider the recurrence:
\[
x_{t+1} = x_t^2 + 1.
\]

This yields:
\[
x_{t+1} - x_t^2 - 1 = 0.
\]

\subsection{Comparison with Circuit Models}

\begin{center}
\begin{tabular}{l|c|c}
Feature & R1CS / PLONK & AIR \\
\hline
View of computation & static circuit & time evolution \\
Basic constraint & $a \cdot b = c$ & $P(s_t, s_{t+1}) = 0$ \\
Structure & graph & trace \\
Best suited for & structured circuits & general programs
\end{tabular}
\end{center}

\subsection{AIR for Machine Learning}

In zkML, AIR can represent inference as a sequence of steps:
\begin{itemize}
\item load inputs
\item apply linear transformations
\item apply nonlinearities
\item produce outputs
\end{itemize}

Each step corresponds to a transition constraint.

\subsection{Key Insight}

\begin{quote}
AIR represents computation as an algebraic dynamical system:
\[
s_{t+1} = F(s_t),
\]
with correctness enforced by polynomial constraints.
\end{quote}

\subsection{Discussion}

\begin{itemize}
\item AIR is more natural for program execution and zkVMs.
\item Circuit-based models are often more compact for structured computations.
\item The choice between AIR and circuit representations reflects a tradeoff between generality and efficiency.
\end{itemize}